\section{Ideas 5 and 6: local windows and telescoping corrections}\label{sec:windows}
\intuition{A seven-position window can have an unusually large swap
count. Its end corrections account for the excess. When the window
slides one position, its new overlap is the old overlap, so the
accounting corrections cancel inside the order.}

\subsection{See one overlap before summing all windows}
Picture eight consecutive positions as
\[
 x\ \big|\ a\ b\ c\ d\ e\ f\ \big|\ y.
\]
The first seven-position window is \(xabcdef\); the next is
\(abcdefy\). Fix the full exterior order in each application.
In the first window, fixing and deleting the first label \(x\)
leaves the central six positions. In the second, fixing and deleting
the last label \(y\) leaves exactly the same six-position family,
with the same ordered exterior. Its correction occurs with a minus
sign in the first inequality and a plus sign in the second.
This is why it cancels. The matching remembers the fixed exterior;
it does not identify unrelated families merely because they look alike.

The weights can also be seen on eight positions. In the following
table all entries have denominator 10:
\begin{center}\small
\begin{tabular}{@{}lrrrrrrr@{}}\toprule
Swap position&1&2&3&4&5&6&7\\\midrule
First window&1&2&2&2&2&1&0\\
Second window&0&1&2&2&2&2&1\\\midrule
Total&1&3&4&4&4&3&1\\\bottomrule
\end{tabular}
\end{center}
There are seven possible swap positions, and the total weight is
\(20/10=2\). The missing weight is \(7-2=5\).
For a longer order the middle positions eventually receive weight
one; the same total deficit of five is confined to the outside ends.
The proof below works uniformly even when the order is short.

\subsection{The complete count}

\begin{proposition}[A global estimate with two outside ends]\label{prop:affine}
There is a fixed constant \(B\), independent of \(n\) and \(K\),
such that for every \(K\in\F_n\) with \(n\geq7\),
\[
 A(K)\leq\frac58n|K|+B|K|,
 \qquad\Sig(K)\geq(3n/8-B)|K|.
\]
\end{proposition}
\begin{proof}
At each window position fix the full ordered prefix and suffix
outside the seven consecutive positions. The remaining labels and
their allowed orders form a conditional family \(J\in\F_7\), unless
empty, by endpoint-slice closure. These conditional families partition
\(K\), so their sizes sum to \(M=|K|\). Different exterior contexts
are retained with their multiplicities even if their local families
are isomorphic. Apply Theorem~\ref{thm:finite} to each context and sum.

The top six-slices of the window in positions \(i,\ldots,i+6\)
match the bottom six-slices of the next window. Both retain positions
\(i+1,\ldots,i+6\) and fix the same ordered exterior. Consequently
all internal \(q\)-terms cancel. Formally this is the elementary sum
\(\sum_{i=0}^{m-1}(Q_i-Q_{i+1})=Q_0-Q_m\), with the sign convention
bottom minus top. The two remaining six-position layers each have
total size \(M\). There are finitely many six-label families and
\(q(\varnothing)=0\), so choose \(Q\) with
\(|q(H)|\leq Q|H|\). The boundary contribution has absolute value
at most \(2QM\). Thus
\[
 \sum_{\text{all window contexts }J}C(J)
 \leq\frac58(n-6)M+2QM.
\]

Every interior swap position has total weight
\((1+2+2+2+2+1)/10=1\) across windows. At all positions the total
weight is between zero and one. There are \(n-1\) swap positions,
but the sum of the weights over the \(n-6\) windows is \(n-6\).
The total missing weight is therefore exactly 5 per order. Since a
swap indicator is at most one,
\[
 A(K)\leq\sum_J C(J)+5M
 \leq\frac58nM+\left(\frac54+2Q\right)M.
\]
Take \(B=5/4+2Q\). Interval fibres give \(\Sig=nM-A\), proving
the second assertion.
\end{proof}
\lean{\path{CDCert.K6.winS_card} and \path{winS_adjAt} identify the
window counts; \path{winS_bot}, \path{winS_top} identify the matching
six-position contexts. \path{window_ineq} applies the finite estimate;
\path{isum_telescope} cancels the internal terms. The final scaled
bound is \path{CDCert.K6.window_bound}:
\[
 8dA\leq5dnM+(10d+2^{66})M.
\]
\path{qD_bound} supplies \(|\texttt{qD}(H)|\leq2^{63}|H|\), so
\(Q=2^{62}/d\) and \(B=5/4+2^{63}/d\) recover exactly this coefficient.
The enormous boundary constant is harmless; the next idea removes it.}
